\chapter{Clifford代数与旋量几何}
\label{ch:19-spin}
\ChapterPlan{周第9章（本地页码待核对）；Figueroa-O'Farrill的Spin讲义及2010习题单；Joyce的特殊Lagrange讲义}{正文证明初稿；8题详解和3幅原创图，尚未独立数学审阅}

第16章给出由主丛表示构造相伴丛的方法；本章的困难是，旋量表示通常不下降为$\mathrm{SO}(n)$的表示。因此既要构造一个更精细的群，也要检查正交标架丛能否提升到这个群。随后，联络的协变导数经Clifford乘法缩并为Dirac算子，其平方把标量曲率带入分析。

除另行声明，$g$正定，流形光滑、定向、无边界，维数$n\geq2$；不默认紧致或完备。采用
\[
 vw+wv=-2g(v,w),\qquad c(v)^*=-c(v),\qquad
 R(X,Y)=[\nabla_X,\nabla_Y]-\nabla_{[X,Y]}.
\]
Clifford代数先在实数上构造，旋量用复Hermite空间，内积第一变量共轭线性。Lorentz旋量、实表示八周期、$\mathrm{Spin}^c$及指标定理不由这些正定论证自动获得。

\section{Clifford代数的构造与表示}
\subsection{用商代数记录度量}
设$(V,g)$为$n$维实内积空间，张量代数$T(V)=\bigoplus_{k\geq0}V^{\otimes k}$取连接乘法。令$I$为所有$v\otimes v+g(v,v)1$生成的双边理想，定义
\[
 \operatorname{Cl}(V,g)=T(V)/I.
\]
把$v+w$的平方减去各自平方，便得$vw+wv=-2g(v,w)$。若线性映射$f:V\to A$进入含幺结合代数且$f(v)^2=-g(v,v)1$，逐个张量相乘唯一延拓到$T(V)$的代数同态，其核含$I$；因此唯一下降为$\operatorname{Cl}(V,g)\to A$。这就是构造的泛性质。还须证明$V$没有在商中意外塌缩。

\Needspace{16\baselineskip}
\begin{theorem}[有序单项式基]\label{thm:clifford-basis}
对任意正交规范基$e_1,\ldots,e_n$，元素
\[
 e_I=e_{i_1}\cdots e_{i_k},\qquad i_1<\cdots<i_k,
 \qquad e_\varnothing=1
\]
构成Clifford代数的基。因此维数为$2^n$，自然映射$V\to\operatorname{Cl}(V,g)$单射。
\end{theorem}
\begin{proof}
用$e_ie_j=-e_je_i$（$i\ne j$）把指标排序，再用$e_i^2=-1$删掉重复指标，说明这些元素生成整个商空间。

为证独立，在$\Lambda^*V$上记$\epsilon_v\alpha=v\wedge\alpha$，$\iota_v$为用$g(v,-)$的内乘。由外代数的交替性和内乘乘法式，
\[
 \epsilon_v\epsilon_w+\epsilon_w\epsilon_v=0,\quad
 \iota_v\iota_w+\iota_w\iota_v=0,\quad
 \iota_v\epsilon_w+\epsilon_w\iota_v=g(v,w)I.
\]
故$C(v)=\epsilon_v-\iota_v$满足所需平方关系，泛性质给表示$C$。有序不同指标作用在$1\in\Lambda^0V$上时，各内乘都遇到不含相同指标的楔积，因此
\[
 C(e_{i_1})\cdots C(e_{i_k})1=e_{i_1}\wedge\cdots\wedge e_{i_k}.
\]
若$\sum_Ia_Ie_I=0$，施加此表示并作用于1，外代数基的独立性迫使全部$a_I=0$。特别地，非零$v$在商中仍非零。
\end{proof}

\subsection{分次、滤过与三个对合}
虽然$v^2$从张量二次降为常数，次数的\emph{奇偶}不变。因此代数有$\mathbb Z/2$分次$\operatorname{Cl}^0\oplus\operatorname{Cl}^1$，乘法将奇偶相加。长度不超过$k$的单项式给滤过$F_k$，伴随分次中的最高项满足$vw=-wv$；由有序基逐次比较，得到
\[
 \operatorname{gr}_F\operatorname{Cl}(V,g)\simeq\Lambda^*V.
\]
反对称化$\frac1{k!}\sum_\sigma\operatorname{sgn}(\sigma)v_{\sigma(1)}\cdots v_{\sigma(k)}$还给出$\Lambda^*V$与Clifford代数的规范\emph{向量空间}同构，但一般不保持乘法：前者$v\wedge v=0$，后者$v^2=-|v|^2$。

定义分次自同构$\alpha(v)=-v$、反序反自同构$\tau(v)=v$、Clifford共轭$\bar a=\alpha(\tau(a))$。它们由保持商关系而存在，且平方均为恒等；$\tau(ab)=\tau(b)\tau(a)$、$\overline{ab}=\bar b\bar a$。这里横线是\emph{实代数反自同构}，不是以后复旋量系数的复共轭。长度$k$的不同正交指标单项式满足
\[
 \alpha(e_I)=(-1)^ke_I,\quad
 \tau(e_I)=(-1)^{k(k-1)/2}e_I,\quad
 \bar e_I=(-1)^{k(k+1)/2}e_I.
\]
例如$\operatorname{Cl}(\mathbb R)=\mathbb C$是\emph{实}代数同构（$e_1\mapsto i$）；二维时$1,e_1,e_2,e_1e_2$满足四元数关系，故$\operatorname{Cl}(\mathbb R^2)\simeq\mathbb H$。这些低维例子不能代替全部实签名及八周期分类。

\subsection{外代数上的复旋量模型}
令$n=2m$，取$S=\Lambda^*\mathbb C^m$，使有序楔积基正交规范。记$\epsilon_j$为楔乘，$\iota_j=\epsilon_j^*$为内乘，定义
\begin{equation}\label{eq:spin-fock}
 c_{2j-1}=\epsilon_j-\iota_j,\qquad
 c_{2j}=i(\epsilon_j+\iota_j).
\end{equation}
外代数恒等式逐项给$c_ac_b+c_bc_a=-2\delta_{ab}I$及$c_a^*=-c_a$，于是$S$为复化Clifford代数的Hermite模。它的维数是$2^m$，不是整个$\Lambda^*\mathbb C^{2m}$的$2^{2m}$。

\Needspace{14\baselineskip}
\begin{proposition}\label{prop:complex-clifford}
复化代数有
\[
 \operatorname{Cl}_{2m}^{\mathbb C}\simeq\operatorname{End}(S),\qquad
 \operatorname{Cl}_{2m+1}^{\mathbb C}\simeq
 \operatorname{End}(S)\oplus\operatorname{End}(S).
\]
偶维有一个不可约复模的同构类，奇维有两个；两者在奇维的偶子代数上同构。
\end{proposition}
\begin{proof}
从$c$能还原
\[
 \epsilon_j=\tfrac12(c_{2j-1}-ic_{2j}),\qquad
 \iota_j=-\tfrac12(c_{2j-1}+ic_{2j}).
\]
令$N_j=\epsilon_j\iota_j$，$P_0=\prod_j(1-N_j)$投影到真空$1$。用适当有序的内乘先把指定楔积基送到1，再接$P_0$和指定的楔乘，得到每一个矩阵单位（其余基向量被投影消去）。因此表示满射；两边维数均为$2^{2m}$，由定理\ref{thm:clifford-basis}知它是同构。

置$\Gamma=i^mc_1\cdots c_{2m}$。把相邻一对相乘可得$\Gamma=\prod_j(1-2N_j)$，即在$\Lambda^k\mathbb C^m$上为$(-1)^k$；它平方为1并与每个$c_a$反交换。奇维另一个生成元可取$c_{2m+1}=+i\Gamma$或$-i\Gamma$。代数中的$z=-i(i^me_1\cdots e_{2m})e_{2m+1}$为中心元且$z^2=1$，两表示分别把它送到$+I,-I$。用$(1\pm z)/2$隔离两个矩阵块，合并表示满射；维数再次相等，得到第二个同构。

最后，矩阵代数的不可约模必是列空间：对一个非零模用对角矩阵单位分解，再用$E_{ab}$相互传送各分量；不可约性迫使每个分量一维。直和代数的中心幂等元在不可约模上仅能有一个非零块。奇维偶子代数的像仍是整个$\operatorname{End}(S)$：其中已有$\Gamma$，再由$c_ac_{2m+1}$与$\Gamma$还原$c_a$。这证明全部断言。
\end{proof}

\subsection{手征与不能下降的表示}
偶维$\Gamma$自伴，给$S=S^+\oplus S^-$，其中$S^\pm$分别为偶/奇外幂，维数均$2^{m-1}$（$m\geq1$）。偶Clifford元素保持它们，向量乘法交换它们；偶子代数的像是$\operatorname{End}(S^+)\oplus\operatorname{End}(S^-)$，因为前述矩阵单位的输入输出同奇偶时其构造属于偶次。

下一节的$\operatorname{Spin}(n)$由偶个单位向量相乘组成，每个$c(v)$在$|v|=1$时为酉矩阵，故其旋量表示酉。$-1$却作用为$-I$，所以此表示不能通过$\operatorname{Spin}(n)\to\mathrm{SO}(n)$下降。一个旋量并非简单地在每张正交标架中放一个复列向量，仍须解决标架过渡提升。


\Needspace{16\baselineskip}
\section{Spin群及其双覆盖}
\subsection{反射生成正交变换}
单位$u\in V$有$u^{-1}=-u$。由Clifford关系直接算得
\[
 uvu=v-2g(u,v)u=:r_u(v).
\]
这是法向为$u$的超平面反射。所有单位向量乘积构成$\operatorname{Pin}(n)$；采用扭伴随$\alpha(a)va^{-1}$时，一个单位$u$的像正是$r_u$。偶个单位向量乘积的子群称$\operatorname{Spin}(n)$，其扭伴随等于普通共轭，记为$\lambda(a)v=ava^{-1}$。

\Needspace{13\baselineskip}
\begin{lemma}[正定反射分解]\label{lem:euclidean-reflections}
每个$T\in\mathrm O(n)$均为至多$n$个超平面反射的乘积；乘积个数的奇偶由$\det T$决定。
\end{lemma}
\begin{proof}
对维数归纳。选单位$e$；若$Te=e$，限制到$e^\perp$后归纳。否则取$u=(Te-e)/|Te-e|$，利用$|Te|=|e|=1$算得$r_uTe=e$。于是$r_uT$固定$e$，在$e^\perp$上由至多$n-1$个反射组成，乘回$r_u$即可。每个反射行列式为$-1$，所以任意这种分解的奇偶一致。此处使用正定性保证非零向量可归一化，未证明不定签名版本。
\end{proof}

\Needspace{17\baselineskip}
\subsection{双覆盖及其微分}
\begin{theorem}\label{thm:spin-cover}
$\operatorname{Spin}(n)$是Lie群，共轭给满射Lie群同态
\[
 1\longrightarrow\{1,-1\}\longrightarrow\operatorname{Spin}(n)
 \xrightarrow{\lambda}\mathrm{SO}(n)\longrightarrow1,
\]
并为二重覆盖。令$L_{ij}e_i=e_j$、$L_{ij}e_j=-e_i$，其余基向量送零，则$\lambda_*(e_ie_j/2)=L_{ij}$。
\end{theorem}
\begin{proof}
每个偶乘积给偶数次反射，像在$\mathrm{SO}(n)$；引理给满射。若偶元素$a$与所有$e_j$交换，按有序基展开：偶长度的$e_I$在$j\notin I$时与$e_j$交换，在$j\in I$时反交换。对固定$j$的交换子涉及的$e_{I\triangle\{j\}}$互不相同，故全部含$j$的系数为零。逐个$j$考察后$a$只能是标量。单位向量满足$\bar u u=1$，乘积仍满足$\bar a a=1$，所以核中的标量只能是$\pm1$。

为说明光滑结构，考虑偶代数中满足$\bar a a=1$且$aV\bar a\subset V$的集合。左乘$a$在有限维代数上有左逆，故可逆且$a^{-1}=\bar a$。共轭保持$v^2=-|v|^2$，所以在$V$上的像正交。偶元素与体积元$e_1\cdots e_n$交换；把一个正交基的像相乘，又等于其行列式乘体积元，所以该正交变换保向。选前述Spin元素$b$具有同样的像，$ab^{-1}$属于偶中心化子且范数为1，故为$\pm1$。这个集合恰是Spin群，由多项式等式闭合；通过忠实左正则表示，它是一般线性群的闭子群，因第15章闭子群定理而为Lie群。

曲线$\cos(t/2)+\sin(t/2)e_ie_j$是两个单位向量的乘积，因而在Spin中。直接取交换子得
\[
 [e_ie_j/2,e_i]=e_j,\qquad [e_ie_j/2,e_j]=-e_i.
\]
故微分像包含整个$\mathfrak{so}(n)$。微分的核中，切向量与$V$交换，前述基论证使其为标量，而对$\bar a a=1$求导使此标量为零。逆函数定理于是给单位元附近局部微分同胚。缩小邻域$U$使$U\cap(-U)=\varnothing$；每个像点的纤维恰为相反的两个元素，故$\lambda^{-1}(\lambda(U))=U\sqcup(-U)$，再用平移即得覆盖。
\end{proof}

对斜对称算子$B$，有特别重要的逆微分公式
\begin{equation}\label{eq:spin-lie-lift}
 \lambda_*^{-1}(B)=\frac14\sum_{i,j}g(Be_i,e_j)e_ie_j.
\end{equation}
这里系数是$g(Be_i,e_j)$，若改用矩阵\emph{行$i$列$j$}的$B_{ij}=g(Be_j,e_i)$，须连同转置号一起改写。

\subsection{二维的半角与三维四元数}
二维偶代数为$\mathbb R\langle1,e_1e_2\rangle$，$(e_1e_2)^2=-1$，故Spin群为单位圆。置$a_\theta=\cos(\theta/2)+\sin(\theta/2)e_1e_2$，共轭使$e_1$旋转到$\cos\theta\,e_1+\sin\theta\,e_2$。提升满足$a_{2\pi}=-1$、$a_{4\pi}=1$。

三维偶代数的$I=e_2e_3,J=e_3e_1,K=e_1e_2$满足$I^2=J^2=K^2=IJK=-1$。每个单位四元数可写为$\cos t+\sin t\,B$，其中单位虚四元数$B$为单位简单双向量：用三维Hodge星把$B$的系数识别为单位法向，再选其正交平面的正向规范基$u,v$，得$B=uv$。此表达是两个单位向量的乘积，所以$\operatorname{Spin}(3)$恰为单位四元数$S^3$，也可用标准二阶复矩阵写成$\mathrm{SU}(2)$。这里没有推断高维Spin群的全局单连通性或全部低维偶然同构。

\begin{figure}[htbp]
\centering\DGDiagram{spin-half-angle}

\caption{平面旋转角$\theta$及其Spin提升的半角。基底转一整周时提升由1到$-1$，两整周才闭合；两圆仅表示参数，不表示旋量是普通平面向量。}
\label{fig:spin-half-angle}
\end{figure}

\subsection{提升的是标架丛，不只是每个过渡矩阵}
设$F_{\mathrm{SO}}(M)$为正向正交标架主丛。一个Spin结构是主$\operatorname{Spin}(n)$丛$Q\to M$，连同覆盖恒等底映射的映射$\Phi:Q\to F_{\mathrm{SO}}(M)$，满足
\[
 \Phi(qa)=\Phi(q)\lambda(a),
\]
且每个纤维上是上述二重覆盖。等价结构要求主丛同构还保持$\Phi$。若只忘掉$\Phi$而比较抽象丛，可能把不同Spin结构错误合并。

在局部标架$s_j=s_i g_{ij}$中，候选提升$h_{ij}:U_i\cap U_j\to\operatorname{Spin}(n)$应满足$\lambda(h_{ij})=g_{ij}$。一个矩阵可提升不保证整个过渡\emph{函数}可单值提升，更不保证三重交集上的乘法相容。

\Needspace{18\baselineskip}
\begin{proposition}[给定覆盖上的符号障碍]\label{prop:spin-cocycle}
取一组有正向正交标架的开覆盖，并假设所有非空有限交集可缩且连通。选$h_{ji}=h_{ij}^{-1}$的光滑提升，令
\[
 z_{ijk}=h_{ij}h_{jk}h_{ki}\in\{1,-1\}.
\]
则$z$为常值层$\mathbb Z/2$的Cech二上闭链，其类不依赖提升选择。此类为零当且仅当这些标架过渡可提升为Spin结构；存在时，等价结构组成$\check H^1(\{U_i\};\mathbb Z/2)$的挠子，即该群作用自由且传递，但未选基准结构时没有指定零点。
\end{proposition}
\begin{proof}
可缩交集上，双覆盖的路径提升使$g_{ij}$有单值提升：从一点选定起值，沿路径提升，收缩路径的同伦提升保证终值与路径无关。覆盖局部逆保证光滑；这一步也可直接用第15章覆盖提升论证。三重乘积在$\lambda$下为1，故是局部常符号；因交集连通而常值。

由$h_{ij}h_{jk}=z_{ijk}h_{ik}$在四重交集计算同一三重乘积的两种结合，得
\[
 z_{ijk}z_{ikl}=z_{jkl}z_{ijl},
\]
其中符号位于中心，正是$\delta z=1$。若$h'_{ij}=\varepsilon_{ij}h_{ij}$，则$z'=z\,\delta\varepsilon$。因此类不变；类为零时选择$\varepsilon$消掉$z$，使$h_{ij}h_{jk}=h_{ik}$，第16章拼接法便构造主Spin丛及$\Phi$。反向由实际Spin丛取局部提升截面，过渡本来就满足该乘法式，障碍类为零。

两组相容提升的比值$\varepsilon_{ij}$为符号一上闭链。保持$\Phi$的局部同构只能改截面为$q_i\eta_i$，$\eta_i\in\{\pm1\}$，因而恰把比值改一个上边界。反之每个符号一上闭链与既有提升相乘仍相容。这给出自由且传递的$\check H^1$作用。
\end{proof}
上闭链、上边界在此均指固定覆盖上交替的符号数据及其乘法，不把它们与第9章的实微分形式混同。\textbf{外引边界：}本书尚未证明每个流形都有上述良覆盖、其Cech群与奇异上同调的比较，以及此障碍等于拓扑$w_2(TM)$的比较。因此常见判据“定向且$w_2=0$当且仅当spin”在采用这些拓扑定理后可用，但不能由第18章实Chern--Weil类直接推出。尤其真实障碍为模2信息，实曲率类一般看不到。

\begin{figure}[htbp]
\centering\DGDiagram{spin-transition-sign}

\caption{三张标架卡的提升相乘可能剩下中心符号$z_{ijk}$。把单个过渡选成另一个提升会改此符号的上边界，但不能任意消掉非零的整体障碍类。}
\label{fig:spin-transition-sign}
\end{figure}

\subsection{可直接构造的例子：平行化与球面}
有全局正向正交标架时，$F_{\mathrm{SO}}\simeq M\times\mathrm{SO}(n)$，取$Q=M\times\operatorname{Spin}(n)$即得一个Spin结构。度量的改变可用第18章正定平方根同构对应正交标架丛，故不改变是否存在Spin结构。

球面不必全局平行化也有Spin结构。在$\mathbb R^{n+1}$取规范基$e_0,e_1,\ldots,e_n$，以
\[
 Q=\operatorname{Spin}(n+1),\qquad p(a)=\lambda(a)e_0\in S^n
\]
为总空间和投影。稳定子恰为嵌入的$\operatorname{Spin}(n)$：其像是固定$e_0$的$\mathrm{SO}(n)$，用$e_0^\perp$内偶数次反射提升；同像元素仅相差$\pm1$，仍在该子群。局部选正交补标架再作局部覆盖提升，给主丛平凡化。映射
\[
 \Phi(a)=(\lambda(a)e_1,\ldots,\lambda(a)e_n)
\]
给$S^n$的正向切标架，其定向采用外向法向在首；它是纤维上的二重覆盖，故确为Spin结构。$S^2$的总空间即$\operatorname{Spin}(3)=S^3$，并不意味着$TS^2$平凡。维数1时同样取$\operatorname{Spin}(1)=\{\pm1\}$，此构造给圆周的连通双覆盖；习题将比较圆周的两个结构。


\Needspace{16\baselineskip}
\section{旋量丛与Dirac算子}
\subsection{相伴旋量丛与提升的联络}
给定Spin结构$(Q,\Phi)$及前述复旋量模$S$，定义$\mathbb S=Q\times_\rho S$，采用第16章关系$[qa,s]=[q,\rho(a)s]$。令$e_i=\Phi(q)e_i^{\mathrm{std}}$，定义
\[
 c(X)[q,s]=[q,c(q^{-1}X)s].
\]
此处$q^{-1}X$意为用标架$\Phi(q)$取坐标，而不是对切向量真的乘一个群逆。由表示的共轭相容性，
\[
 \rho(a)c(v)\rho(a)^{-1}=c(\lambda(a)v),
\]
所以定义与$q$无关。偶维$\Gamma$也不依赖正向标架，给$\mathbb S^+\oplus\mathbb S^-$；反转定向会改变手征标记。

Levi--Civita主联络$\omega$拉回到$Q$后，再用$\lambda_*^{-1}$识别李代数，得到$\widetilde\omega=\lambda_*^{-1}\Phi^*\omega$。它满足主联络的再现及等变公理，因$\lambda$为群同态且微分同构。相伴联络写成
\begin{equation}\label{eq:spin-connection}
 \nabla_X^{\mathbb S}\psi=X(\psi)+A^{\mathbb S}(X)\psi,
 \qquad A^{\mathbb S}(X)=\frac14\sum_{i,j}
 g(\nabla_Xe_i,e_j)c_i c_j.
\end{equation}
相邻提升$q_j=q_i h_{ij}$中，局部系数按$\rho(h_{ij})^{-1}$变换，联络势按第16章公式变换，故上式确实拼接。

\Needspace{13\baselineskip}
\begin{proposition}\label{prop:spin-compatible}
旋量联络保持Hermite内积，并满足
\[
 \nabla_X^{\mathbb S}(c(Y)\psi)=c(\nabla_XY)\psi+c(Y)\nabla_X^{\mathbb S}\psi.
\]
其曲率为
\begin{equation}\label{eq:spin-curvature}
 R^{\mathbb S}(X,Y)=\frac14\sum_{i,j}
 g(R(X,Y)e_i,e_j)c_i c_j.
\end{equation}
\end{proposition}
\begin{proof}
对$i\ne j$，$(c_ic_j)^*=c_jc_i=-c_ic_j$，实斜系数使$A^{\mathbb S}$斜Hermite，因此联络度量相容。直接算
\[
 [c_ic_j,c_k]=2(\delta_{ik}c_j-\delta_{jk}c_i),
\]
连同\eqref{eq:spin-lie-lift}说明$[A^{\mathbb S}(X),c_k]=c(\nabla_Xe_k)$。对$Y=\sum Y^ke_k$用乘法式即得相容性。曲率可从主丛同态及表示拉回，也可在$\dd A^{\mathbb S}+A^{\mathbb S}\wedge A^{\mathbb S}$中用同一交换子计算；因$\lambda_*^{-1}$为李代数同态，结果正是\eqref{eq:spin-curvature}。
\end{proof}

\subsection{Dirac算子、主符号与形式自伴}
定义几何Dirac算子
\begin{equation}\label{eq:geometric-dirac}
 D\psi=\sum_{i=1}^nc(e_i)\nabla^{\mathbb S}_{e_i}\psi.
\end{equation}
若换为$e'_i=\sum_j e_jO_{ji}$，两处$O$相乘求和给$\sum_iO_{ji}O_{ki}=\delta_{jk}$，所以定义与标架无关。偶维$D$交换手征：$D:\Gamma(\mathbb S^\pm)\to\Gamma(\mathbb S^\mp)$，因为$\Gamma$平行且与向量乘法反交换。

\Needspace{14\baselineskip}
\begin{proposition}\label{prop:dirac-symbol-adjoint}
有$[D,f]=c(\operatorname{grad}f)$。以$\partial_j\mapsto i\xi_j$定义主符号，则
\[
 \sigma_D(x,\xi)=i c(\xi^\sharp),\qquad
 \sigma_D(x,\xi)^2=|\xi|^2I,
\]
故$D$椭圆。对紧支光滑旋量$\psi,\chi$，有$\int\langle D\psi,\chi\rangle=\int\langle\psi,D\chi\rangle$。
\end{proposition}
\begin{proof}
展开$D(f\psi)$，只有导数作用于$f$的项未相消，正是$\sum_i(e_if)c_i$。符号把它再乘$i$，由Clifford平方得正的$|\xi|^2$。不用$i$的实代数符号约定也可用，但其平方会是$-|\xi|^2$，不能混写。

在一点取协变导数为零的正交标架，对复值向量场$W=\sum_i\langle\psi,c_i\chi\rangle e_i$，相容性和$c_i^*=-c_i$给
\[
 \operatorname{div}W=-\langle D\psi,\chi\rangle+\langle\psi,D\chi\rangle.
\]
这是标量恒等式，逐点成立。对实虚部应用Stokes，紧支使边界项为零，即得结论。
\end{proof}
形式自伴只讨论积分分部；非紧时的算子闭包、定义域、本质自伴、紧致时的谱完备和Fredholm性质需要分析论证，不由上述命题自动获得。若有边界，还须写出边界项及边界条件，本章不处理APS边界问题。

\subsection{同样的局部公式，不同的全局符号}
平坦$T^n=\mathbb R^n/\mathbb Z^n$上，任选$\varepsilon_j\in\{0,\tfrac12\}$，令
\[
 \psi(x+\ell)=e^{2\pi i\varepsilon\cdot\ell}\psi(x),\qquad \ell\in\mathbb Z^n.
\]
由$(x,a)\sim(x+\ell,e^{2\pi i\varepsilon\cdot\ell}a)$构造Spin主丛（相位这里只是中心$\pm1$），再取相伴丛便得这些条件。这$2^n$个结构互不等价且穷尽：一个Spin结构在固定全局标架的逆像是二重覆盖；拉回可缩的$\mathbb R^n$后，用路径及同伦提升将其平凡化。格平移在两层上给乘法特征$\mathbb Z^n\to\{\pm1\}$，恰由$n$个基圈符号决定，保持$\Phi$的同构不改变这些符号。

固定平行标架下$D=\sum c_j\partial_j$，但允许的旋量不同。每个模式
\[
 \psi(x)=e^{2\pi i(k+\varepsilon)\cdot x}s
\]
被$D$送为同一指数乘$2\pi i\,c(k+\varepsilon)s$；非零频率的平方特征值为$4\pi^2|k+\varepsilon|^2$。对$n\geq2$非零频率的正负特征值均出现，重数相等，可用与频率正交的Clifford乘法交换两特征空间。

此处先给显式模式，不以一组模式的写出代替一般谱完备证明。下一节的能量恒等式会直接说明：平坦环面调和旋量必平行，所以只有全周期结构有常旋量核。局部曲率均为零，却可以有不同的全局核，这正是Spin结构选择不可忽略的例子。

\subsection{Dirac平方的曲率项}
定义联络粗Laplace算子
\[
 \nabla^*\nabla\psi=-\sum_i\left(
 \nabla^{\mathbb S}_{e_i}\nabla^{\mathbb S}_{e_i}\psi
 -\nabla^{\mathbb S}_{\nabla_{e_i}e_i}\psi\right).
\]
它在函数平凡线丛上是第11章$-\Delta_g$，在旋量上仍有非负的紧支积分$\int\langle\nabla^*\nabla\psi,\psi\rangle=\int|\nabla\psi|^2$；对每个局部标架用度量相容和散度，再用单位分解即可证明此积分式。

\Needspace{15\baselineskip}
\begin{theorem}[Schr\"odinger--Lichnerowicz公式]\label{thm:lichnerowicz}
对正定Riemannian Spin流形的未扭旋量丛，有算子恒等式
\begin{equation}\label{eq:lichnerowicz}
 D^2=\nabla^*\nabla+\tfrac14\mathrm{Scal}\,I.
\end{equation}
\end{theorem}
\begin{proof}
在任意点$p$选局部正向正交标架使$\nabla e_i(p)=0$；无挠使$[e_i,e_j](p)=0$。相容性、$c_i^2=-1$及不同指标的反交换给
\[
 D^2=\nabla^*\nabla+\sum_{i<j}c_ic_jR^{\mathbb S}(e_i,e_j)
 \quad\text{在$p$处}.
\]
记$R_{ijkl}=g(R(e_i,e_j)e_k,e_l)$，代入\eqref{eq:spin-curvature}，曲率项为
\[
 \mathcal C=\frac18\sum_{i,j,k,l}R_{ijkl}c_ic_jc_kc_l.
\]
下面不用“标准缩并”省掉关键号。第11章给两对反对称、对换对称$R_{ijkl}=R_{klij}$及第一Bianchi。

若两对指标只共一个，交换$(i,j)$与$(k,l)$时系数不变，而对应两个双Clifford乘积反交换，所以这样的项成对抵消。四指标互异时，对每个$i<j<k<l$合并全部排列，剩下
\[
 (R_{ijkl}-R_{ikjl}+R_{iljk})c_ic_jc_kc_l=0,
\]
等于零正是第一Bianchi和对换对称。若仅有两个不同指标$i<j$，四个非零排列用$(c_ic_j)^2=-1$合并，给$R_{ijji}/2$。因此
\[
 \mathcal C=\tfrac12\sum_{i<j}R_{ijji}
 =\tfrac14\sum_{i,j}R_{ijji}=\tfrac14\mathrm{Scal}.
\]
这里$R_{ijji}$是正截面曲率的分子，与本书球面为正的约定一致。$p$任意，故算子恒等式全局成立。
\end{proof}

\Needspace{13\baselineskip}
\subsection{调和旋量消失及其适用边界}
\begin{corollary}\label{cor:spin-vanishing}
设$M$紧致、连通、无边界且spin，$\mathrm{Scal}\geq0$。若$D\psi=0$，则$\psi$平行；若某点$\mathrm{Scal}>0$，则$\psi=0$。
\end{corollary}
\begin{proof}
形式自伴及上式给
\begin{equation}\label{eq:spin-energy}
 \|D\psi\|_{L^2}^2=\|\nabla\psi\|_{L^2}^2
 +\frac14\int_M\mathrm{Scal}|\psi|^2\,\mathrm{vol}.
\end{equation}
左端为零，右边两项非负，故平行。度量相容使其长度常数；若正曲率点附近积分仍为零，旋量在该处为零，连通性使常长度为零。若不连通，只能逐分支使用此结论。
\end{proof}
单位$n$球面有$\mathrm{Scal}=n(n-1)>0$（$n\geq2$），前节又显式构造了Spin结构，所以没有非零调和旋量。平坦环面则只有全周期结构允许非零平行旋量，核维数为$2^{\lfloor n/2\rfloor}$，其余符号结构的核为零。此核计算仅用能量式及平行方程，不依赖Fourier谱完备。

非紧空间上不能直接把无穷远的边界项丢掉；有边界时更不能省略边界条件。正标量曲率不能仅凭本章结论排除\emph{所有}Spin流形上的度量：要从拓扑指标强迫核非零，仍需第20章计划中的Fredholm及指标理论。

\subsection{扭丛会多出曲率}
若$(E,h,\nabla^E)$为复Hermite丛及酉联络，在$\mathbb S\otimes E$上用张量积联络和$c\otimes I$定义$D_E$。曲率相加，以上推导逐字保留到最后一步，给
\begin{equation}\label{eq:twisted-lichnerowicz}
 D_E^2=\nabla^*\nabla+\tfrac14\mathrm{Scal}
 +\sum_{i<j}c_ic_j\otimes F^E(e_i,e_j).
\end{equation}
最后一项一般没有非负性，不能只查标量曲率就重复未扭的消失结论。二维平面上取$A^E=iBx\dd y$，有$F^E=iB\dd x\wedge\dd y$、$\Gamma=ic_1c_2$，故
\[
 D_E^2=-\partial_x^2-(\partial_y+iBx)^2+B\Gamma.
\]
系数及号可直接通过两个协变导数的交换子核对；本例不声称已发展规范场的谱理论。


\Needspace{16\baselineskip}
\section{校准与几何应用}
\subsection{闭形式怎样给体积下界}
几何中的校准（calibration）不是调节参数，而是一个闭形式给出的逐切平面体积下界。设$\varphi$为Riemannian流形上的实$k$形式，要求对每个有向正交规范$k$标架有$\varphi(v_1,\ldots,v_k)\leq1$；反转定向也在要求中，故等价于绝对值不超过1。若再有$\dd\varphi=0$，称为校准。定向$k$维子流形$N$满足$\varphi|_N=\mathrm{vol}_N$时，称被其校准。

\Needspace{14\baselineskip}
\begin{theorem}[校准的体积比较]\label{thm:calibration-minimizing}
设$N,N'$为紧致定向无边界$k$维子流形，且存在紧致定向光滑带边$(k+1)$维流形$W$及光滑$F:W\to M$，使边界映射给$N'$和反向的$N$。若$N$被$\varphi$校准，则$\operatorname{Vol}(N)\leq\operatorname{Vol}(N')$。此外，校准子流形对紧支变分为极小子流形。
\end{theorem}
\begin{proof}
第8章Stokes给$\int_{N'}\varphi-\int_N\varphi=\int_W\dd F^*\varphi=0$。逐点的校准不等式给$\int_{N'}\varphi\leq\operatorname{Vol}(N')$，而$N$上等于体积，得到比较。

为得局部极小性，取支撑在一张子流形坐标卡内部的变分，选包含支撑的闭坐标球域。Cartan公式给$\frac{\dd}{\dd t}\int F_t^*\varphi=\int F_t^*\dd\iota_V\varphi=0$，因为变分场在边界附近为零，Stokes的边界项消失。只比较该紧区域体积，即使$N$非紧也成立。$t$的正负两侧体积都不小于原值，导数为零；第13章第一变分用任意局部法向场遂得平均曲率向量处处为零。分解一般紧支变分后亦成立。
\end{proof}
通常的“同调类内最小”还可用光滑链版本的Stokes表述；这里直接证明给定光滑配边的版本及固定边界的局部变分，不预用一般奇异同调比较、积分流存在、紧性或奇点正则性。

\subsection{Kahler形式的Wirtinger不等式}
设$(M^{2m},g,J,\omega)$为Kahler，沿用$\omega(X,Y)=g(JX,Y)$。则$\varphi_k=\omega^k/k!$为校准，且其校准的有向$2k$平面恰为复定向的$J$不变平面。证明如下：在实$2k$平面$W$上，$A=\operatorname{pr}_WJ|_W$为斜自伴且算子范数不超过1。实斜矩阵的正交分块标准形给正交基，使$\omega|_W$由$k$个系数$\lambda_j$的面积块组成，$|\lambda_j|\leq1$；于是
\[
 \left|\frac{\omega^k}{k!}(v_1,\ldots,v_{2k})\right|
 =\prod_j|\lambda_j|\leq1.
\]
等号使每个块$|\lambda_j|=1$，故$|\operatorname{pr}_WJv|=|v|$；正交投影的等号迫使$Jv\in W$。反向$J$不变且取复定向时，由第17章体积式得正的等号。$\dd\omega=0$使$\varphi_k$闭，因此复子流形在相应光滑配边范围内体积最小。若只是Hermite，逐点不等式仍成立，但缺少闭性不能作Stokes比较。

\subsection{平坦特殊Lagrange与相位}
在$\mathbb C^m$上取Euclid度量、$\omega_0=\sum_j\dd x_j\wedge\dd y_j$、$\Omega_0=\dd z_1\wedge\cdots\wedge\dd z_m$，则$\operatorname{Re}\Omega_0$是校准。确实，把有向实规范基$v_1,\ldots,v_m$作为复矩阵$Z$的列，每列Hermite长度为1，复Hadamard不等式给
\[
 \operatorname{Re}\det Z\leq|\det Z|\leq1.
\]
Hadamard可由复Gram--Schmidt证明：行列式绝对值是各列依次投影到前列正交补后的长度之积，每项至多原长度；等号当且仅当各列复正交。由于各列已实正交，复正交等价于$\omega_0|_W=0$。再要求$\det Z=1$，即$\operatorname{Im}\Omega_0|_W=0$且实部取正，才是当前定向下的相位0校准；只要求虚部为零可能给相反定向。

\begin{figure}[htbp]
\centering\DGDiagram{calibration-phase}

\caption{在一族对角Lagrange平面中，$\Omega_0(v_1,\ldots,v_m)=e^{i\sum\theta_j}$。相位0给正向校准，实部非正的分支不能仅因虚部为零就混入；图只画该复数的相位圆。}
\label{fig:calibration-phase}
\end{figure}

对实函数$f:\mathbb R^m\to\mathbb R$，图像$x\mapsto x+i\operatorname{grad}f$的切矩阵为$I+iH$，$H=\operatorname{Hess}f$对称。拉回$\omega_0$为零，而
\begin{equation}\label{eq:special-lagrangian-graph}
 \Omega_0|_{\Gamma_f}=\det(I+iH)\,\dd x_1\cdots\dd x_m,
 \qquad \mathrm{vol}_{\Gamma_f}=\sqrt{\det(I+H^2)}\,\dd x_1\cdots\dd x_m.
\end{equation}
对角化$H$给$\sqrt{\det(I+H^2)}=|\det(I+iH)|$。因此采用投影到$x$空间的图像定向时，相位0方程必须是
\[
 \operatorname{Im}\det(I+iH)=0,\qquad \operatorname{Re}\det(I+iH)>0.
\]
这些图被校准，因而对紧支光滑比较体积最小。本章不证明其一般非线性PDE解的存在或模空间理论。

\subsection{旋量与校准之间还差哪些论证}
旋量的作用在于给平行张量的一种来源。若$\psi$平行，对任意$k$，张量
\[
 \beta(X_1,\ldots,X_k)=\operatorname{Alt}
 \langle\psi,c(X_1)\cdots c(X_k)\psi\rangle
\]
其中$\operatorname{Alt}$采用$1/k!$归一化。由联络相容性逐项求导为零；取实部或虚部仍为平行实形式，Levi--Civita无挠使其外微分为零。然而它可能为零，或其最大值未归一化为1，等号平面也尚未确定。要从特定旋量得到$G_2$、$\operatorname{Spin}(7)$校准，必须另证表示恒等式、稳定子和逐平面不等式。本章仅建立上述平行闭性与Kahler/平坦特殊Lagrange实例，不把特殊holonomy分类、Killing旋量分类、Calabi--Yau存在性及McLean变形理论算作已证。


\Needspace{18\baselineskip}
\section{旋量几何习题及详细解答}
19.1--19.3分别从Figueroa-O'Farrill的2010年Spin课程习题单3的题3.3、3.5、3.6改述、限制或扩展，19.4改述习题单4的题4.1；两份习题单各只有一个PDF页，均印刷/PDF1。19.5--19.8为原创补充。讲义正文是2017年5月18日修订版；来源符号与本书显式定义分别核对，不照搬可能不同的签名或对合记号。

\Needspace{12\baselineskip}
\subsection*{习题19.1：限制到Spin后何时仍不可约}
改述习题单3题3.3，取正定情形。证明不可约Clifford模限制到Pin仍不可约，不可约偶Clifford模限制到Spin仍不可约。为什么这不使偶维的整个$S=S^+\oplus S^-$成为不可约Spin模？
\begin{proof}[解答]
每个规范基向量为单位向量，任意Clifford有序基单项式都是Pin元素，故Pin的线性张成是整个Clifford代数。一个对Pin不变的线性子空间，对其线性张成也不变，因此不可约性保留。同理，偶有序基单项式属于Spin，故Spin张成偶子代数，得到第二个结论；这对实模或复化后的复模均成立。

偶维的全$S$只对\emph{整个}Clifford代数不可约；对偶子代数，它已分成两个非零的手征块，所以第二个结论的假设不适用于全$S$。将它分别用于$S^+$、$S^-$，才得到两个不可约半旋量表示。原题不是“所有Clifford不可约模限制到Spin都不可约”。
\end{proof}

\Needspace{12\baselineskip}
\subsection*{习题19.2：三个对合的号与逆元}
改述习题单3题3.5，采用本章$\alpha,\tau,\bar{\ }$的定义。求它们在$k$个不同正交生成元乘积上的号，验证$\overline{ab}=\bar b\bar a$，并计算单位向量乘积的逆。
\begin{proof}[解答]
分次对每个生成元给一个负号，故为$(-1)^k$。反序需$k(k-1)/2$次相邻交换，故$\tau$的号为$(-1)^{k(k-1)/2}$；合成为$(-1)^{k(k+1)/2}$。在$k=0,1,2,3$模4时，三列号依次是
\[
 \begin{array}{c|rrrr}
 k\bmod4&0&1&2&3\\\hline
 \alpha&+&-&+&-\\
 \tau&+&+&-&-\\
 \bar{\ }&+&-&-&+
 \end{array}
\]
反序对一般乘积先交换两个因子的先后顺序，分次再逐因子作用，便得共轭的反乘法式。对单位$u$有$u^{-1}=-u=\bar u$，故$a=u_1\cdots u_r$满足
\[
 a^{-1}=(-u_r)\cdots(-u_1)=\bar a.
\]
这些号取决于明确给出的定义；不可只因不同资料都称某记号“共轭”就把符号表照抄。旋量矩阵中的Hermite伴随则还含系数复共轭，是另一种运算。
\end{proof}

\Needspace{12\baselineskip}
\subsection*{习题19.3：Spin3的一个矩阵实现}
限制并扩展习题单3题3.6，只处理$\operatorname{Spin}(3)$。把$I,J,K$送到$-i\sigma_1,-i\sigma_2,-i\sigma_3$（$\sigma_j$为第18章Pauli矩阵），证明单位四元数映成$\mathrm{SU}(2)$，说明覆盖核。
\begin{proof}[解答]
Pauli乘法给$(-i\sigma_j)^2=-I$及$(-i\sigma_1)(-i\sigma_2)=-i\sigma_3$，循环式同样成立，故为实代数同态。四个像实线性独立，四元数$q=a+bI+cJ+dK$的像满足$U^*U=(a^2+b^2+c^2+d^2)I$，行列式也等于这个实数。单位四元数因此进入SU2。

反向SU2矩阵的两列酉正交且行列式1，迫使其形如$\left(\begin{smallmatrix}z&w\\-\bar w&\bar z\end{smallmatrix}\right)$，$|z|^2+|w|^2=1$；它能唯一写成上述四个基矩阵的实组合。因此映射为群同构。共轭在三维虚四元数上为正交保向变换；按前文三维Hodge识别，它与Spin覆盖的切向量表示对应。与所有虚四元数交换的$q$为实数，单位条件迫使$q=\pm1$，故核正是两点。此题未证明Spin4以上的同构表。
\end{proof}

\Needspace{12\baselineskip}
\subsection*{习题19.4：平坦二维环面的四个符号}
改述习题单4题4.1，使用环面平行标架作基准。以$b=(b_1,b_2)\in(\mathbb Z/2)^2$记录绕两基圈的符号$(-1)^{b_j}$。求$A\in\mathrm{SL}(2,\mathbb Z)$所给换格映射的拉回作用，以及四个结构的轨道。
\begin{proof}[解答]
符号是格上的特征$\chi_b(\ell)=(-1)^{b^{\mathsf T}\ell}$。换格$x\mapsto Ax$把闭路$\ell$送到$A\ell$，故拉回为$\chi_b\circ A$，即
\[
 b\longmapsto A^{\mathsf T}b\pmod2.
\]
导数$A$为常矩阵；比较拉回度量与原度量可用常的正定平方根，正交标架的比较同样为常，故不会额外引入随闭路改变的符号。这也说明此线性表达使用了平行标架的基准。

生成矩阵$S=\left(\begin{smallmatrix}0&-1\\1&0\end{smallmatrix}\right)$、$T=\left(\begin{smallmatrix}1&1\\0&1\end{smallmatrix}\right)$分别给$(b_2,b_1)$、$(b_1,b_1+b_2)$。零向量固定，其余三点由这些变换互相连通，故有一个单点轨道和一个三点轨道。这里只写这些显式换格映射的作用；不在本题证明任意环面微分同胚均由此代表或高属曲面的映射类群分类。
\end{proof}

\Needspace{12\baselineskip}
\subsection*{习题19.5：二维旋量与半角}
原创补充。在$\Lambda^*\mathbb C$的基$(1,u)$中写出$c_1,c_2,\Gamma$及$\rho(a_\theta)$。计算$\theta=2\pi$时的值，并说明半旋量是否是SO2的单值表示。
\begin{proof}[解答]
楔乘把1送到$u$，内乘把$u$送到1，故
\[
 c_1=\begin{pmatrix}0&-1\\1&0\end{pmatrix},\quad
 c_2=\begin{pmatrix}0&i\\i&0\end{pmatrix},\quad
 \Gamma=ic_1c_2=\begin{pmatrix}1&0\\0&-1\end{pmatrix}.
\]
于是$\rho(a_\theta)=\operatorname{diag}(e^{-i\theta/2},e^{i\theta/2})$。在$2\pi$处两对角元都是$-1$，$4\pi$处才回到1。因此两个一维半旋量均为Spin2的表示，却不是以$\theta\bmod2\pi$为参数的SO2单值表示。向量共轭作用仍在$2\pi$返回恒等，两种表示不要混淆。
\end{proof}

\Needspace{12\baselineskip}
\subsection*{习题19.6：圆周的边界Spin结构与Dirac模式}
原创补充。长度$L$的圆周采用$c(e)=i$。分别写出周期和反周期结构的Dirac模式及核。哪一个由圆盘的Spin结构按外向法向在首诱导？
\begin{proof}[解答]
用弧长$x$取平行切标架，则$D=i\partial_x$。周期条件$\psi(x+L)=\psi(x)$允许模式$e^{2\pi ikx/L}$，其特征值$-2\pi k/L$；反周期条件允许$e^{2\pi i(k+1/2)x/L}$，特征值$-2\pi(k+1/2)/L$。从$D\psi=0$直接得常数，故前者核复维1，后者为零，无需先用Fourier完备性。

圆盘内固定Spin平凡化。沿边界，外向法向和正向切向量组成的正向标架转一整周；Spin2提升从1走到$-1$，故诱导的Spin1双覆盖是连通的，旋量相对于边界平行切标架呈反周期。这是反周期结构，不能因为圆盘自身可平凡化就把其边界的旋转标架当作常标架。
\end{proof}

\Needspace{12\baselineskip}
\subsection*{习题19.7：扭Dirac平方的二维号}
原创补充。Euclid平面上令$\nabla_x=\partial_x$、$\nabla_y=\partial_y+iBx$，$B$为实常数。直接展开$D_E=c_1\nabla_x+c_2\nabla_y$的平方，比较两个手征分量，并解释为何不能丢掉扭曲率项。
\begin{proof}[解答]
对角项为$-\nabla_x^2-\nabla_y^2$，交叉项为$c_1c_2[\nabla_x,\nabla_y]=iBc_1c_2=B\Gamma$。因此在$\Gamma=\pm1$的分量上，零阶修正分别为$\pm B$。取$B>0$，其中一个修正为负，故粗Laplace加上未扭标量曲率的非负性论证不能直接复用。这里验证的是局部微分算子恒等式；平面非紧，既未给L2定义域，也未据此宣布完整Landau能级或紧致扭核的维数。
\end{proof}

\Needspace{12\baselineskip}
\subsection*{习题19.8：图像方程的相位分支}
原创补充。对$\Gamma_f\subset\mathbb C^m$取投影图像定向。证明二维时相位0特殊Lagrange方程等价于$\partial_1^2f+\partial_2^2f=0$；三维时写出方程，并用$\operatorname{Hess}f=\sqrt3 I$说明只查虚部为零的错误。
\begin{proof}[解答]
二维实对称$H$有$\det(I+iH)=1-\det H+i\operatorname{tr}H$。虚部零等价于$\Delta_g f=\operatorname{tr}H=0$；此时两特征值为$\lambda,-\lambda$，实部$1+\lambda^2>0$，所以二维不需再排除另一分支。

三维展开给
\[
 \det(I+iH)=1-\sigma_2(H)+i\bigl(\operatorname{tr}H-\det H\bigr),
\]
因此必须同时有$\operatorname{tr}H=\det H$和$1-\sigma_2(H)>0$。取$f=(\sqrt3/2)|x|^2$，三个特征值均$\sqrt3$，则$\det(I+iH)=(1+i\sqrt3)^3=-8$。虚部虽零，图像定向下实部却为负，故是相位$\pi$而非相位0；反转定向后才被$\operatorname{Re}\Omega_0$校准。此反例说明相位条件具有实际内容，不是无关的记号选择。
\end{proof}

